\paragraph{What the same-prime quotient does retain.}
The preceding cross-prime rank loss must not be confused
with a loss of the entire zeta function. There is a general
elementary complement. Let \(F\) be a profinite group,
let \(K\triangleleft F\) be closed, fix a prime \(r\),
and put \(N=\ker(K\to K^{(r)})\). For every open
subgroup \(I\leq F\) containing \(N\), the natural map
induces an isomorphism
\[
                         I^{(r)}\simeq (I/N)^{(r)}.
                                                        \tag{AG.39.15}
\]
Here the superscript denotes the maximal pro-\(r\)
quotient. In particular, the assertion is not restricted
to open subgroups having an \(r\)-group Galois core.

To prove it, take an open normal subgroup \(J\triangleleft I\)
such that \(I/J\) is an \(r\)-group, and set
\(I_0=\operatorname{core}_F(I)\).
The group \(J\cap I_0\) is normal in \(I_0\), and
\(I_0/(J\cap I_0)\) embeds in \(I/J\).
Let \(M=\operatorname{core}_F(J\cap I_0)\).
Then \(I_0/M\) embeds in a finite product of
\(r\)-groups, so it is an \(r\)-group. Since
\(N\triangleleft F\) and \(N\subset I\), we have
\(N\subset I_0\); consequently the image of \(K\)
in \(F/I_0\) is an \(r\)-group. Its image in
\(F/M\) is an extension of an \(r\)-group by a
subgroup of \(I_0/M\), and is therefore an
\(r\)-group as well. The defining property of \(N\)
now gives \(N\subset M\subset J\). Thus every
finite \(r\)-quotient of \(I\) kills \(N\),
which proves (AG.39.15).

Apply this to the geometric punctured-curve group and
its cuspidal kernel. For every finite cover visible in
the maximal cuspidally pro-\(r\) quotient, quotienting
by the puncture inertia and then taking the maximal
pro-\(r\) quotient recovers the \emph{full} proper
pro-\(r\) group of its smooth compactification.
Indeed, formation of a maximal pro-\(r\) quotient
commutes with quotienting by the closed normal subgroup
generated by the specified inertia groups. This applies
in particular to the ramified cover \(D\) in
(AG.39.11), even though its other-prime cohomology has
been lost.

If the curves and cover are defined over a finite field
of characteristic different from \(r\), this comparison
is compatible with the arithmetic Frobenius action.
It therefore retains \(H^1(D_{\bar k},\mathbb Z_r)\)
with Frobenius, and hence the exact integral Weil
polynomial and zeta function of \(D\). Equality of the
\(r\)-adic characteristic polynomials is equality of
their integer polynomials, since
\(\mathbb Z[T]\hookrightarrow\mathbb Z_r[T]\).
Thus the obstruction in (AG.39.11) concerns the full
proper prime-to-\(p\) group needed by a subsequent
different-prime reconstruction theorem; it does not
prevent transfer of these zeta functions. No inference
from that transfer to point separation is made here.
In particular, the Weil polynomial does not by itself
specify the integral \(p\)-primary rational torsion
exponent, or supply R1-compatible point lifts on these
ramified covers.
